\chapter{摘\quad 要}

一篇文章读起来累不累，取决于两件事：写了什么，以及排成了什么样子。
已有的可读性公式大多只看前者——句子多长、词有多难——而把版面当作不变的背景。
对于动辄上万字的中文长文本，这种做法会漏掉一个重要的变量：
同样的文字，排在 40 字一行的屏幕上和 20 字一行的书页上，读者付出的努力并不相同。

本文提出一种同时考虑\textbf{语言难度}与\textbf{版面负荷}的中文长文本可读性评估方法。
主要工作如下：
（1）构建了包含 1\,200 篇文章、36\,000 条阅读记录的中文长文本阅读数据集，
记录了每位读者在不同版式下的阅读时间与理解得分；
（2）提出版面负荷指标，用行长、行距与段落密度刻画「排版让眼睛多做了多少事」；
（3）设计了融合语言特征与版面特征的评估模型，在数据集上与人工评分的相关系数达到
$0.81$，比只用语言特征的基线提高 $0.12$。

实验表明，版面因素可以解释可读性差异中约五分之一的部分。
这一结论为中文排版规范、阅读类产品的版式设计提供了可量化的依据。

\vspace{1em}
\noindent{\heiti 关键词：}可读性；中文排版；长文本；阅读负荷；评估模型

\chapter*{Abstract}
\addcontentsline{toc}{chapter}{Abstract}
\markboth{Abstract}{Abstract}

Whether a text is tiring to read depends on what it says and on how it is set.
Classical readability formulas look only at the former. For long-form Chinese text,
which may run to tens of thousands of characters, this ignores an important variable:
the same words cost the reader different effort on a 40-character screen line and on
a 20-character book line.

This thesis proposes a readability assessment method for long-form Chinese text that
accounts for both \emph{linguistic difficulty} and \emph{layout load}. We build a
dataset of 1,200 articles and 36,000 reading sessions, define a layout-load index from
line length, leading and paragraph density, and design a model that combines linguistic
and layout features. The model correlates with human ratings at $r = 0.81$, an
improvement of $0.12$ over a language-only baseline.

\vspace{1em}
\noindent\textbf{Keywords:} readability; Chinese typography; long-form text;
reading load; assessment model
